\documentclass[10pt,a4paper]{amsart}
\usepackage[margin=19mm]{geometry}
\usepackage{amsmath,amssymb,mathtools}
\usepackage[scr=rsfs,cal=euler]{mathalfa}
\usepackage{xfrac}
\usepackage[unicode,hidelinks,bookmarks=false]{hyperref}
\newcommand{\ie}{i.e.\ }
\newcommand{\cf}{cf.\ }
\newcommand{\resp}{respectively\ }
\newcommand{\Subsection}{\subsection}
\providecommand{\nobreakdash}{\nobreak\hbox{-}}
\setlength{\emergencystretch}{2em}
\multlinegap0pt
\usepackage{mathtools}
\usepackage{amssymb}
\usepackage{stackengine}
\usepackage{bm}
\usepackage{xcolor}
\usepackage[all]{xy}
\usepackage{enumerate}
\usepackage[shortlabels]{enumitem}
\usepackage{tikz}
\usepackage{caption}
\newtheorem{rem}{Remark}
\newtheorem{prop}{Proposition}
\newcommand{\CC}{\mathbb{C}}
\newcommand{\TT}{\mathbb{T}}
\newcommand{\ZZ}{\mathbb{Z}}
\newcommand{\ga}{\gamma}
\title{Geometric realisation of hypergeometric local systems}
\author{Asem Abdelraouf, Giulia Gugiatti}
\date{Source: arXiv:2602.14204v1 (CC BY 4.0)}
\begin{document}
\maketitle

\noindent\emph{Source and licence.} Geometric realisation of hypergeometric local systems. Version arXiv:2602.14204v1 (CC BY 4.0), distributed by arXiv under the Creative Commons Attribution 4.0 International licence, http://creativecommons.org/licenses/by/4.0/, as recorded in the arXiv licence metadata for this exact version and shown on the abstract page https://arxiv.org/abs/2602.14204v1. Full licence notice: \texttt{licenses/arxiv-2602.14204-CC-BY-4.0.txt}.\par\smallskip

\noindent\textit{Editorial scope.} Counted unit: the article's prop prop:u-result (source0.tex lines 1385--1396). The article's complete original proof of it is delivered with it (source0.tex lines 1398--1484, one \texttt{proof} environment of 87 lines). No mathematics is written, repaired or re-derived: the statement and the proof are the article's own lines, in the article's own notation, and no retained line is substituted or re-flowed.  Retained uncounted as the source-local context the proof needs: lines 1289--1302; lines 1304--1310; lines 1312--1322.  Nothing was omitted from any retained span, so the package carries no omission record.  No retained line contains a citation, so the wrapper carries no bibliography and no bibliography entry.  No retained line contains a figure environment, an image-inclusion command or drawing code, so no image is delivered and the package declares no asset dependency.  Editorial formatting only, in addition: an AMS \texttt{amsart} preamble that restates the article's own declarations and macros used by the retained lines, namely the environments \texttt{rem}, \texttt{prop} on separate counters, together with the standard packages those lines need.  The retained environments print in this document as Remark 1, Proposition 1 in order of appearance, whereas the article prints the same items with the numbering of its own full document.  The complete native source of the article as distributed in the arXiv e-print for this exact version is bundled unchanged as \texttt{sources/194695/source0.tex}.
\begin{rem}
\label{rem:gamma-relations}
The vector $\gamma$ spans the lattice of integral relations among the vectors 
    \begin{equation}
        \label{eq:KT}
    \begin{pmatrix}
1  \\
m_{j} 
\end{pmatrix} \in \ZZ^{l-1}\quad 
     j=1, \dots, l
     \end{equation}

This implies that the vectors $m_j$ are all distinct. Equivalently, for all $t\in \CC^\times$,  $f_t$ is the sum of $l$ distinct monomial terms. 
\end{rem}

\begin{rem}
\label{rem:kernel}
The first $l-1$ rows of $\mathcal{A}$ span a saturated sublattice of $\ZZ^l$ of rank $l-1$. 
Equivalently, the maximal minors of the $(l-1)\times l$ submatrix of $A$ defined by the first $l-1$ rows are relatively prime. In particular, the vectors 
\eqref{eq:KT}
    span $\ZZ^{l-1}$. This is easily seen to be equivalent to the fact that, for any $k \in \{1, \dots, l\}$, the vectors $m_j-m_k$, $j=1, \dots, l$, span $\ZZ^d$. 
\end{rem}

\begin{rem}
\label{rem:mj-polytope}
The polytope $\Delta$ is full-dimensional and simplicial\footnote{A polytope is called simplicial if all its proper faces are simplices.}. 
Indeed, since for any fixed $k$ the vectors $m_j-m_{k}$ generate $\ZZ^d$, $\Delta$ has dimension $d$. Then, either $\Delta$ has $d+1$ vertices, i.e., it is a simplex, or it has $d+2$ vertices. In the second case, any facet of $\Delta$ must be a simplex: if a facet $F \subset \Delta$ had $d+1$ vertices, say $m_2, \dots, m_l$, then  for all $j \neq 1$,  $\langle m_j, u\rangle =b$ for some $u \in \ZZ^d, b \in \ZZ$, while $\langle m_1, u\rangle =a>b$; on the other hand, 
\[
0=\left\langle \sum_{j=1}^l \gamma_j m_j, u \right\rangle =\sum_{j=1}^l \gamma_j \langle m_j, u \rangle=  \gamma_1a+\sum_{j=2}^l \gamma_j b   = \ \gamma_1 (a -b)
\] 
implies that $a=b$ since $\ga_1 \neq 0$. 

It is simple to see that $\Delta$ is a simplex if and only if either $\gamma$ or $-\gamma\coloneqq (-\gamma_1, \dots, -\gamma_l)$  has only one negative entry. More precisely, if $\gamma_k$ is the only negative/positive entry of $\gamma$, then $m_k$ lies in the interior of $\Delta$. 
\end{rem}

\begin{prop}
\label{prop:u-result} 

Let $\mathcal{F}=\sum_{j=1}^l u_j x^{m_j} \in \CC[x_1^{\pm 1}, \dots, x_d^{\pm 1}, u_1^{\pm 1}, \dots, u_1^{\pm 1}]$. Let $\mathcal{Z}=Z_\mathcal{F}$ be the zero locus of $\mathcal{F}$ in $\TT^d \times (\CC^\times)^l$. 
For $u=(u_1, \dots, u_l) \in (\CC^\times)^l$ fixed, let $\mathcal{F}_u=\mathcal{F}(\cdot,u)$, and let $\mathcal{Z}_u=Z_{\mathcal{F}_u}$ be the zero of $\mathcal{F}_u$ in $\TT^d$. Then:   
    \begin{enumerate}
    \item  If $\prod_{j=1}^l u_j^{\gamma_j} \neq \Gamma$, then $\mathcal{Z}_u$ is smooth. 
    \item  If $\prod_{j=1}^l u_j^{\gamma_j} = \Gamma$, then $\mathcal{Z}_u$ has a unique ODP.  
    \item  For all $u \in (\CC^\times)^l$, $\mathcal{Z}_u$ is quasi $\Delta$-regular. In particular, if $\prod_{j=1}^l u_j^{\gamma_j} \neq \Gamma$, then $\mathcal{Z}_u$ is $\Delta$-regular. 
    \item $\mathcal{Z}$ is smooth.
    \end{enumerate}
\end{prop}

\begin{proof}[Proof of Proposition \ref{prop:u-result}] 
Fix $u \in (\CC^\times)^l$.
For $i=1, \dots, d$, we have
\begin{equation}
 x_i \frac{\partial  \mathcal{F}_u }{\partial x_i} = \sum_{j=1}^l m_{ij} u_j x^{m_j}
\end{equation}
The variety $\mathcal{Z}_u$ is singular if and only there is a point  $x=(x_1, \dots, x_d) \in \TT^d$ that satisfies the equations
\begin{equation}
\label{eq:singularity}
\sum_{j=1}^l u_j x^{m_j}=0,  \quad 
  \sum_{j=1}^l m_{ij}  u_j x^{m_{j}}=0 \ \text{ for } i =1,\dots,d 
\end{equation} 
which we can rewrite as
\begin{equation*}
\label{eq:1m-system}
\begin{pmatrix}
1 & \dots & 1\\
m_1 & \dots & m_l
\end{pmatrix} \begin{pmatrix}
u_1 x^{m_1}\\
\vdots\\
u_l x^{m_l}
\end{pmatrix} = 0
\end{equation*}
Then, by Remark \ref{rem:gamma-relations}, $x \in \TT^d$ satisfies \eqref{eq:singularity} if and only if 
    \begin{equation}
    \label{eq:y-gamma}
    (u_1 x^{m_1}, \dots, u_l x^{m_l}) =(c \gamma_1,  \dots, c \gamma_l)\end{equation} for some constant $c \in \CC^\times$. 
If such $x$ exists, then 
\[
\prod_{j=1}^l u_j^{\gamma_j}= \prod_{j=1}^l \left( u_j x^{m_j} \right)^{\gamma_j}= \prod_{j=1}^l\left(  c\gamma_j\right)^{\gamma_j}= c^{\sum_{j=0}^l \gamma_j} \cdot \Gamma=\Gamma
\]
which proves $(1)$. 

Assume now that $u \in (\CC^\times)^l$ is such that $\prod_{j=1}^l u_j^{\gamma_j}=\Gamma$. 
Then there is a unique $x \in \TT^d$ satisfying \eqref{eq:y-gamma}. 
Indeed, if \eqref{eq:y-gamma} holds, then, for $j=1, \dots, l$, 
\[ x^{m_j-m_1}=\frac{u_1 \gamma_j}{\gamma_1 u_j}
\]
Since the vectors $m_j-m_1$, $j=2, \dots, l$ generate generate $\ZZ^d$ (see Remark \ref{rem:kernel}), these identities determine $x=(x_1, \dots, x_d)$ uniquely (note that $c=u_1x^{m_1}/\gamma_1$ is also uniquely determined).
This shows that, if $\prod_{j=1}^l u_j^{\gamma_j}=\Gamma$, then $\mathcal{Z}_u$ has a unique singular point. To show $(2)$, it remains to show that the singular point is an ODP. 

Up to changing the variables $x_i$ by some constants, we may assume that $u=\gamma$. Then the singular point of $\mathcal{Z}_u$ is $x=(1, \dots, 1) \in \TT^d$. The Hessian of $\mathcal{F}_u$ at $x=(1, \dots, 1)$ is the matrix $H$ with entries 
\begin{equation}\label{eq:hessian}
H_{i,j}= \sum_{h=1}^l \gamma_h m_{i,h} m_{j,h}  x^{m_h}\bigg|_{x=1} = \sum_{h=1}^l \gamma_h m_{i,h} m_{j,h}.
\end{equation}
Equivalently, letting $M$ be the $d \times l$ submatrix of ${A}$ with entries $m_{ij}$ and $\mathrm{diag}(\gamma)$  be the diagonal matrix with diagonal entries $\gamma_1, \dots, \gamma_l$, we have that $H=M \cdot \mathrm{diag}(\gamma) \cdot M^T$.
Moreover, 
the product ${A} \cdot \mathrm{diag}(\gamma)  \cdot {A}^T$ is of the form: 
\begin{equation}
\label{eq:prod-H}
{A} \cdot \mathrm{diag}(\gamma)  \cdot {A}^T =\begin{bmatrix}
0 & 0 & 1 \\
0 & H & * \\
1 & * & *
\end{bmatrix} 
\end{equation}
Since $\det({A})=\pm 1$, one has $\det({A} \cdot \mathrm{diag}(\gamma)  \cdot {A}^T)=\prod_{j=1}^l \gamma_j$. Combining this with \eqref{eq:prod-H} yields  $\prod_{j=1}^l \gamma_j=-\det(H)$,  thus $H$ is invertible. This shows that $x$ is an ODP.

To prove $(3)$, we verify that, for each face $F$ of $\Delta$ of dimension $r \in (0,d)$, the zero locus of ${\mathcal{F}_u}_{|F}$ in $\TT^d$ is smooth. 
For any face $F\subset \Delta$, the zero locus of ${\mathcal{F}_u}_{|F}$ in $\TT^d$ is singular if and only if there exists $x \in \TT^d$ such that 
\begin{equation}
\label{eq:singularity-face}
   \sum_{m_j \in F } u_j x^{m_{j}}=0, \quad  
    \sum_{m_j \in F }   m_{ij} u_j x^{m_j}=0  \ \text{ for } i=1,\dots,d
\end{equation}
Let $F\subset \Delta$ be a face of dimension $r$. Then $F$ is a $r$-simplex, i.e., it has $r+1$ vertices $m_{j_1}, \dots, m_{j_{r+1}}$, and does not contain any other vertex $m_j$ (see Remark \ref{rem:mj-polytope}). 
Then, we can rewrite  
\eqref{eq:singularity-face} as 
\begin{equation}
\label{eq:ls-face}
\begin{pmatrix}
1 & \cdots  & 1\\
m_{j_1} & \cdots & \ m_{j_{r+1}}
\end{pmatrix}\begin{pmatrix}
u_1 x^{m_1}\\
\vdots\\
u_l x^{m_l}
\end{pmatrix} = 0
\end{equation}
Since any maximal minor of the $(l-1)\times l$ submatrix of ${A}$ defined by its first $l-1$ rows is nonzero (see Remark \ref{rem:kernel}) and $r+1<l-1$,  \eqref{eq:ls-face} does not admit any non-trivial solution.  

Statement $(4)$ follows immediately from the fact that, for all $h=1, \dots, l$,
\[u_h \frac{\partial \mathcal{F}} {\partial u_h}= u_h x^{m_h}\]
and that, for each $k$ fixed, the vectors $m_j-m_k$ generate $\ZZ^d$.

\end{proof}

\end{document}
